\section{Results}
\subsection{Hydrolysis of the Polyimide}
The hydrolysis of the polyimide follows a standard mechanism for hydrolysis of an amide to a carboxylic acid using sodium hydroxide (NaOH). These show the hydroxy group attacking the carbonyl which then breaks the C-N bond and gives the carboxylic acid free.

During the formation of the hydrolysed version, the polymer sheets were converted into a light orange suspension which was filtered to give a pale orange precipitate.

\subsubsection{Structural Analysis}
Structural determination was conducted via Infrared spectroscopy. A comparison was made between the attempted hydrolysis of the sheets and a purchased form (polyamic acid). 

\subsection{Cross-linking of the Polyimide}
Cross-linking attempts follow a mechanism as proposed by ( put reference here), where ethylene carbonate is used to form an ester with the carboxylic acids on neighbouring polymer chains, thereby joining them together with a "cross-link". This was performed with three different salts ( TEAB, $ZnCl_2$, $FeCl_3$) as catalysts. 

During the cross-linking reaction with TEAB, an orange-brown suspension was formed that when washed, filtered and freeze dried gave a pale orange powder. With $ZnCl_2$, the liquid was a dark orange clear liquid that, when water was added, gave a single clump of grey polymeric residue. This was freeze dried to remove the excess water and then analysed. The cross-linking reaction with the $FeCl_3$, gave a dark orange/brown gel, which gave a grey gum-like texture in small pieces when water was added. This was also frozen and analysed. 

\subsubsection{Infrared Spectrum analysis}
